\section[Cote 118 : produits stables sans asphéricité]{Cote 118 : produits stables sans asphéricité\protect\verifie{Folder118}}\label{sec:118}

La cote 118, \emph{Cohomology properties of maps in (Cat)}, quarante pages que l'inventaire date \q{[à partir de 1983]}, énonce page~32, avec un point d'interrogation, un lemme de conditions équivalentes pour une petite catégorie $A$. Les pages~33 et~34 établissent que (i) \q{$A$ totalement asphérique} équivaut à (ii) \q{$\mathcal W_A$ stable par produits finis} jointe à l'asphéricité de $A$. La page~35 demande si une catégorie connexe qui satisfait (ii) est asphérique, et répond \q{Non} : pour $A$ un groupoïde, $\mathcal W_A$ est formé des seuls isomorphismes, donc stable par produits finis (\lot{118}{2}{35}). La connexité, l'hypothèse $G \neq 1$ et la raison donnée pour \q{formé des seuls iso} sont de la lecture.

\noindent\emph{Conventions.} Pour $A = BG$, les préfaisceaux sur $A$ sont les $G$-ensembles, et $A_{/X}$ est le groupoïde d'action $X /\!/ G$. Le modèle est tronqué en degré~$1$ : un foncteur entre groupoïdes est une équivalence faible s'il induit une bijection sur les composantes connexes et des isomorphismes sur les groupes d'automorphismes, et un groupoïde est asphérique s'il est connexe à groupes d'automorphismes triviaux. Ces caractérisations, classiques, servent de définitions et ne sont pas démontrées. Ainsi $f : X \to Y$ est dans $\mathcal W_A$ si $f$ induit une bijection sur les orbites et si $\mathrm{Stab}(f(x)) = \mathrm{Stab}(x)$ pour tout $x$.

\begin{contrexemple}[Grothendieck, page 35, la réponse \q{Non}\piste{118-groupoid-products-not-totally-aspherical}]
Soit $G$ un groupe non trivial et $A = BG$. Alors $\mathcal W_A$ est formé des bijections équivariantes, donc stable par produits de toute famille, finie ou non. Pourtant $A = A_{/e}$ n'est pas asphérique, ni le produit $a \times a$ du représentable $a$ par lui-même.
\end{contrexemple}

\begin{proof}
Si $f \in \mathcal W_A$ et $f(x) = f(x')$, la bijectivité sur les orbites donne $h$ avec $hx = x'$ ; alors $h \in \mathrm{Stab}(f(x)) = \mathrm{Stab}(x)$ et $x' = x$. Tout $y$ s'écrit $g f(x) = f(gx)$ : $f$ est bijective. La réciproque est immédiate, et un produit de bijections est une bijection. Le groupe d'automorphismes de l'unique objet de $A_{/e}$ est $G \neq 1$. Enfin $a$ est le $G$-ensemble $G$, et dans $G \times G$ les points $(1, 1)$ et $(1, g)$, $g \neq 1$, ne sont pas dans la même orbite : $A_{/a \times a}$ n'est pas connexe.
\end{proof}

\begin{remarque}
L'énoncé tient, et pour les produits infinis aussi, que la page laisse en suspens dans le cas général. Il confirme ce que la piste relève chez Cisinski~\cite{cisinski2006}, proposition~4.3.2 : pour $A = BG$, la condition (g), stabilité des $\mathcal W$-équivalences de préfaisceaux par produits finis, est satisfaite, tandis que la condition (a), asphéricité des $a \times b$, ne l'est pas. Dans le modèle tronqué (\texttt{answer\_page\_35}), la première n'entraîne donc pas la seconde ; savoir si l'énoncé publié restreint (g) reste la question de la piste.
\end{remarque}

Ne sont pas démontrés ici : l'identification des équivalences faibles de groupoïdes à celles du modèle tronqué, les équivalences (i) $\Leftrightarrow$ (ii) $+$ \q{$A$ asphérique} et (iv) $\Leftrightarrow$ asphéricité des $a \times b$ des pages~33 et~34, et le corollaire~2 de la page~31.
